%%%% ABSTRACT

\begin{abstractutfpr}%% Ambiente abstractutfpr
Home advantage is one of the traditional pillars of football culture, yet the
tactical, logistical and technological transformations of the modern era raise
questions about its actual magnitude and about the mechanisms that sustain it.
This work quantitatively analyses the influence of playing at home in the
Premier League, covering eleven complete seasons, from 2015/2016 to 2025/2026,
across a dataset of 4,180 matches. The research follows the CRISP-DM methodology
and uses the Python language for data collection, cleaning, statistical analysis
and modelling. In addition to descriptive and inferential analysis, a bivariate
Poisson model with a home advantage term is estimated, which quantifies the
advantage while controlling for opponent strength, and predictive models are
evaluated under temporal validation. The results confirm the existence of the
advantage with a high degree of confidence: home teams win 58.09 percent of
decided matches and obtain 0.37 more points per game, corresponding to a 21.7
percent increase in expected goals. The associated effect size, however, is
small, and the advantage manifests itself through a greater volume of
opportunities rather than through superior finishing efficiency. Three findings
contradict current expectations. First, home advantage increases with team
quality: clubs finishing among the top six obtain a larger differential than
those finishing among the bottom six, in a monotonic gradient repeated across
all analysed metrics. Second, the disciplinary bias favouring the home team
disappears in matches without spectators already in the 2019/2020 season, while
the advantage in points remains unchanged, dissociating refereeing bias from the
central mechanism of the phenomenon. Third, attributing the decline of the
advantage in the 2020/2021 season to the absence of crowds is not supported when
the contrast is constructed within a single season. In the predictive stage, the
statistical model outperforms the evaluated machine learning models, and the
calibration analysis demonstrates why accuracy is insufficient as an evaluation
criterion in probabilistic problems.
\end{abstractutfpr}
